% GENERATED FILE. Edit canonical lessons, then run npm run generate:books.
% canonical-source-hash: fnv1a64:98c4cb6aad53f6c5
% canonical-lessons: FR-C51-peau, FR-C51-ongle, FR-C51-menton, FR-C51-gorge, FR-C51-joue

\chapter{Skin, Nail, Chin, Throat, Cheek}
\label{ch:fr-skin-nail-chin-throat-cheek}
\hlchaptermodality{51}

\begin{chapteropening}
\textbf{By the end of this chapter:} \emph{I can name skin, a nail (finger or toe), the chin, the throat, a cheek.}

Complete the last lesson of chapter 51: I can name skin, a nail (finger or toe), the chin, the throat, a cheek.
\end{chapteropening}

\section[la peau]{la peau — skin}

\label{lesson:FR-C51-peau}



\begin{quote}
Before the new one: say the French for the face, then the French for blood.
\end{quote}

\subsection*{You'll want to know: la peau}
\textbf{la peau} — \textquotedblleft{}skin\textquotedblright{}.

\begin{cousinweb}
\textbf{peau} $\leftarrow$ Old French \emph{pel} $\leftarrow$ Latin \emph{pellis}, \textquotedblleft{}skin, hide\textquotedblright{} — inherited. (EtymDB 2.1)
\end{cousinweb}

\begin{grammarlens}[title={What you've built}]
One more word for this chapter.
\end{grammarlens}

\subsection*{Your turn}
\begin{itemize}
\raggedright
  \item \emph{Say it:} \emph{la peau}
  \item \emph{Say it:} \emph{la peau}, once more
  \item \emph{Say it:} say \emph{la peau}
  \item \emph{Recall:} say the French for the face, then the French for blood, then say \emph{la peau} again
\end{itemize}

\begin{tcolorbox}[breakable,colback=teal!4,colframe=teal!35!black,title={Before you move on}]
What does la peau mean? (\textquotedblleft{}Skin\textquotedblright{}.) Say it once more.
\end{tcolorbox}
\section[l'ongle]{l'ongle — a nail (finger or toe)}

\label{lesson:FR-C51-ongle}



\begin{quote}
Before the new one: say the French for blood, then the French for skin.
\end{quote}

\subsection*{You'll want to know: l'ongle}
\textbf{l'ongle} — \textquotedblleft{}a nail (finger or toe)\textquotedblright{}.

\begin{cousinweb}
\textbf{ongle} $\leftarrow$ Latin \emph{ungula}, \textquotedblleft{}a claw, a hoof\textquotedblright{}, from \emph{unguis}, \textquotedblleft{}nail\textquotedblright{} (Lewis \& Short) — inherited. EtymDB 2.1 traces Italian \emph{unghia} to the same \emph{ungula}.
\end{cousinweb}

\begin{grammarlens}[title={What you've built}]
One more word for this chapter.
\end{grammarlens}

\subsection*{Your turn}
\begin{itemize}
\raggedright
  \item \emph{Say it:} \emph{l'ongle}
  \item \emph{Say it:} \emph{l'ongle}, once more
  \item \emph{Say it:} say \emph{l'ongle}
  \item \emph{Recall:} say the French for blood, then the French for skin, then say \emph{l'ongle} again
\end{itemize}

\begin{tcolorbox}[breakable,colback=teal!4,colframe=teal!35!black,title={Before you move on}]
What does l'ongle mean? (\textquotedblleft{}A nail (finger or toe)\textquotedblright{}.) Say it once more.
\end{tcolorbox}
\section[le menton]{le menton — the chin}

\label{lesson:FR-C51-menton}



\begin{quote}
Before the new one: say the French for skin, then the French for a nail.
\end{quote}

\subsection*{You'll want to know: le menton}
\textbf{le menton} — \textquotedblleft{}the chin\textquotedblright{}.

\begin{cousinweb}
\textbf{menton} $\leftarrow$ Vulgar Latin \emph{*mento}, from Latin \emph{mentum}, \textquotedblleft{}chin\textquotedblright{} — inherited. (EtymDB 2.1)
\end{cousinweb}

\begin{grammarlens}[title={What you've built}]
One more word for this chapter.
\end{grammarlens}

\subsection*{Your turn}
\begin{itemize}
\raggedright
  \item \emph{Say it:} \emph{le menton}
  \item \emph{Say it:} \emph{le menton}, once more
  \item \emph{Say it:} say \emph{le menton}
  \item \emph{Recall:} say the French for skin, then the French for a nail, then say \emph{le menton} again
\end{itemize}

\begin{tcolorbox}[breakable,colback=teal!4,colframe=teal!35!black,title={Before you move on}]
What does le menton mean? (\textquotedblleft{}The chin\textquotedblright{}.) Say it once more.
\end{tcolorbox}
\section[la gorge]{la gorge — the throat}

\label{lesson:FR-C51-gorge}



\begin{quote}
Before the new one: say the French for a nail, then the French for the chin.
\end{quote}

\subsection*{You'll want to know: la gorge}
\textbf{la gorge} — \textquotedblleft{}the throat\textquotedblright{}.

\begin{cousinweb}
\textbf{gorge} $\leftarrow$ Late Latin \emph{gurga}, from Latin \emph{gurges}, \textquotedblleft{}a whirlpool, an abyss\textquotedblright{} — inherited. (EtymDB 2.1)
\end{cousinweb}

\begin{grammarlens}[title={What you've built}]
One more word for this chapter.
\end{grammarlens}

\subsection*{Your turn}
\begin{itemize}
\raggedright
  \item \emph{Say it:} \emph{la gorge}
  \item \emph{Say it:} \emph{la gorge}, once more
  \item \emph{Say it:} say \emph{la gorge}
  \item \emph{Recall:} say the French for a nail, then the French for the chin, then say \emph{la gorge} again
\end{itemize}

\begin{tcolorbox}[breakable,colback=teal!4,colframe=teal!35!black,title={Before you move on}]
What does la gorge mean? (\textquotedblleft{}The throat\textquotedblright{}.) Say it once more.
\end{tcolorbox}
\section[la joue]{la joue — a cheek}

\label{lesson:FR-C51-joue}



\begin{quote}
Before the new one: say the French for the chin, then the French for the throat.
\end{quote}

\subsection*{You'll want to know: la joue}
\textbf{la joue} — \textquotedblleft{}a cheek\textquotedblright{}.

\begin{cousinweb}
\textbf{joue} $\leftarrow$ Old French \emph{jode} $\leftarrow$ Vulgar Latin \emph{*gauta}, \textquotedblleft{}cheek\textquotedblright{} — inherited (EtymDB 2.1). Where \emph{*gauta} came from is uncertain; Gaulish is the usual guess.
\end{cousinweb}

\begin{grammarlens}[title={What you've built}]
One more word for this chapter.
\end{grammarlens}

\subsection*{Your turn}
\begin{itemize}
\raggedright
  \item \emph{Say it:} \emph{la joue}
  \item \emph{Say it:} \emph{la joue}, once more
  \item \emph{Say it:} say \emph{la joue}
  \item \emph{Recall:} say the French for the chin, then the French for the throat, then say \emph{la joue} again
\end{itemize}

\begin{tcolorbox}[breakable,colback=teal!4,colframe=teal!35!black,title={Before you move on}]
What does la joue mean? (\textquotedblleft{}A cheek\textquotedblright{}.) Say it once more.
\end{tcolorbox}
